% Introduction
\section{Introduction}
\label{sec:introduction}

The emergence of large language models (LLMs) capable of generating, understanding,
and reasoning about code has sparked intense interest in multi-agent systems for
software engineering automation \citep{li2023chatdev, hong2023metagpt}. These systems
promise to decompose complex software projects into subtasks, assign them to specialized
agents, and coordinate their execution toward a unified goal. Early results demonstrate
that multi-agent collaboration can produce functioning software artifacts, from simple
games to multi-file applications \citep{qian2023chatdev}.

However, current multi-agent approaches face three fundamental challenges:

\paragraph{Coordination Overhead.} Systems like ChatDev \citep{qian2023chatdev} and
MetaGPT \citep{hong2023metagpt} rely on sequential role-playing or message-passing
between agents, creating bottlenecks that limit parallelism. When Agent A must wait
for Agent B's output before proceeding, potential speedups from parallel execution
remain unrealized.

\paragraph{Quality Control Gaps.} Most multi-agent systems lack structured validation
between planning and execution phases. Agents execute plans without rigorous assessment
of plan quality, leading to wasted computation when fundamental flaws are discovered
late. Similarly, post-execution quality assessment is often informal or absent.

\paragraph{Limited Human Oversight.} Fully autonomous multi-agent systems exclude
human judgment from the loop, missing opportunities to leverage human expertise for
task decomposition, ambiguity resolution, and quality assessment. When agents
encounter edge cases or make suboptimal decisions, there is no mechanism for
human correction.

\subsection{SPOQ: Our Approach}

We introduce \textbf{SPOQ} (Specialist Orchestrated Queuing), a methodology that
addresses these challenges through three integrated innovations:

\begin{enumerate}
    \item \textbf{Wave-Based Topological Dispatch:} We model task dependencies as a
    directed acyclic graph (DAG), a structure where arrows show which tasks must
    complete before others can begin, with no circular dependencies (think of it
    as a flowchart where all arrows point forward). We then compute execution
    \textit{waves} (groups of independent tasks) via topological sort. Tasks within
    the same wave execute in parallel, while waves execute sequentially to respect
    dependencies. This maximizes parallelism without coordination overhead.

    \item \textbf{Dual Validation Gates:} We apply structured validation (quality
    checkpoints with scored metrics) at two points: \textit{before} execution
    (planning validation with 10 metrics) and \textit{after} execution (code
    validation with 10 metrics). Each gate enforces a 95\% threshold, catching
    quality issues when they are cheapest to fix.

    \item \textbf{Human-as-an-Agent (HaaA):} A human specialist participates alongside
    AI agents, not as a passive observer, but as an active collaborator who
    decomposes tasks, validates plans, and can be consulted during execution.
    This bidirectional integration treats the human as a high-value agent within
    the system.
\end{enumerate}

\subsection{Contributions}

This paper makes the following contributions:

\begin{itemize}
    \item A formal framework for wave-based multi-agent orchestration that
    computes parallel execution waves from task dependency
    graphs~(Section~\ref{sec:methodology})

    \item A three-tier agent hierarchy (Opus/Sonnet/Haiku) that optimizes
    cost-quality tradeoffs by matching model capability to task
    complexity~(Section~\ref{subsec:agent-hierarchy})

    \item The Human-as-an-Agent (HaaA) paradigm for structured task
    decomposition through bidirectional human-AI
    collaboration~(Section~\ref{subsec:haaa})

    \item A dual validation system with explicit metrics that scores both
    planning quality and code quality against quantified
    thresholds~(Section~\ref{sec:validation})

    \item A controlled benchmark suite testing four research questions:
    scheduling efficiency, planning quality, validation effectiveness, and
    human-AI collaboration. Wave dispatch reaches the critical-path lower
    bound (ratio 1.03--1.11, up to 14.3$\times$ speedup) on unbounded
    synthetic DAGs and delivers a stable 1.4$\times$ speedup matching the
    hardware ceiling on a 2-slot local backend;
    structured planning lifts coverage to 99.75 and parallelism to 75.25;
    dual validation reduces defects from 0.34 to 0.20 per task; and
    Human-as-Agent planning further reduces defects to 0.03 per task with
    99.75\% test pass rate~(Section~\ref{sec:evaluation})

    \item A cross-provider replication of all four experiments against a
    locally hosted open-weights model (Qwen3.6-35B-A3B served via
    \texttt{llama.cpp}). The direction and significance of the SPOQ gains
    are preserved across provider families, supporting the claim that the
    improvements stem from orchestration rather than from a specific
    model's capabilities~(Section~\ref{sec:evaluation})

    \item A longitudinal deployment study across 17 repositories, 8{,}589
    commits, 1{,}822 completed tasks, and 13{,}866 executed tests
    (99.87\% aggregate pass rate) demonstrating ecological validity,
    third-party adoption, and operational viability at
    scale~(Section~\ref{sec:evaluation})
\end{itemize}

\begin{glossarybox}[Key Terms at a Glance]
\begin{description}
\item[Epic:] A high-level goal decomposed into atomic tasks, each with explicit dependencies, acceptance criteria, and time estimates.
\item[Wave:] A group of tasks sharing no mutual dependencies, enabling simultaneous execution by multiple agents within a single phase.
\item[DAG:] Directed Acyclic Graph, a structure where arrows represent prerequisite relationships between tasks, with no circular dependencies.
\item[Critical Path:] The longest sequential chain of dependent tasks through the DAG; determines the minimum possible project duration.
\item[Validation Gate:] A scored quality checkpoint where work must exceed defined metric thresholds before the pipeline advances to the next stage.
\item[HaaA (Human-as-an-Agent):] A bidirectional collaboration paradigm where humans participate alongside AI agents as active contributors.
\item[PERT Estimates:] Three-point time estimates capturing optimistic, realistic, and pessimistic durations to quantify scheduling uncertainty.
\end{description}
\end{glossarybox}

\newpage
